\subsection{Ordinary class actions and the remaining prime-class transvection}

\begin{statement}{E.91}{Proved}{All vertical ordinary actions, the original class-group types, and an exact ninth-genus refinement}
Use the literal fields, roots, prime labels and norm maps of E.86:
\[
 L_d=\mathbb Q(\zeta,\beta_d),\qquad
 B_d=L_d(\xi),\qquad
 \xi^3=\zeta,\quad \beta_{d+1}^3=\beta_d,\quad
 \beta_d^{3^d}=17\zeta .
 \tag{E91.1}
\]
Write $\mathcal C_d=\operatorname{Cl}(B_d)\{3\}$, let $R_d\subset \mathcal C_d$
be the subgroup generated by its six primes above $3,17$, and put
$D_d=\operatorname{Cl}(L_d)\{3\}$. The three-primary ideal
classes are written additively. The actual vertical generator
$\sigma_d(\beta_{d+1})=\zeta\beta_{d+1}$ acts trivially on
the \emph{entire ordinary} group $\mathcal C_{d+1}$, for every $d\ge1$.
Consequently the adjacent ordinary norm and extension maps satisfy
\[
 \begin{gathered}
 N_d:\mathcal C_{d+1}\twoheadrightarrow \mathcal C_d,\qquad
 j_dN_d=3,\qquad
 \ker N_d\subset \mathcal C_{d+1}[3],\\
 \ker j_d=N_d\bigl(\mathcal C_{d+1}[3]\bigr).
 \end{gathered}
 \tag{E91.2}
\]
These are maps of ordinary ideal classes; their further quotients
by $R_d$ are the specified S-class maps of E.86.

Let $q_d=[\mathfrak q_{3,d}]$, and write $\alpha_d$ for the
exponent called $a_d$ in E69.5, so that
$\operatorname{ord}(q_d)=3^{\alpha_d}$. For every $d\ge2$,
\[
 D_d\simeq C_{3^{\alpha_d+1}}\oplus C_3,
 \qquad \langle q_d\rangle=3D_d,
 \qquad D_d/3D_d\xrightarrow{\sim}\operatorname{Cl}_S(L_d)\{3\}.
 \tag{E91.3}
\]
The last arrow is the actual removal-of-S-primes map; its basis is
the pair of specified ideal classes in E86.4. In particular those
two ordinary ideal classes generate $D_d$. Adjacent norms on the
$D_d$ are surjective and send $q_{d+1}$ to $q_d$. Thus the whole
ordinary tower, as well as its specified prime, stabilizes under
norm once the proved finite exponent $h$ of E.69 stabilizes:
\[
 \varprojlim_N D_d\simeq C_{3^{h+1}}\oplus C_3,
 \qquad (q_d)_d\text{ generates its subgroup of triples}.
 \tag{E91.4}
\]
The isomorphism is a choice of a cyclic basis in the finite stable
group. Its map to every level is the actual iterated ideal norm.

The next unevaluated ordinary class group has a small, explicit
residual comparison. On $B_3$ put
\[
 g(\xi)=\xi,\quad g(\beta_3)=\xi\beta_3,
 \qquad h(\xi)=\zeta\xi,\quad h(\beta_3)=\beta_3.
 \tag{E91.5}
\]
Let $\mathfrak p_3$ be the E.86 prime above $73$ with
$\xi=55$ and $\beta_3=22$. Then
\[
 \begin{gathered}
 W=(g-1)\mathcal C_3=\langle(g-1)[\mathfrak p_3]_3\rangle,
 \qquad 3W=0,\quad |W|\le3,\quad h|_W=1,\\
 0\longrightarrow W\longrightarrow\mathcal C_3
   \xrightarrow{\ \mathrm{Art}\ } A_9\longrightarrow0,
 \qquad A_9=C_9^2\oplus C_3 .
 \end{gathered}
 \tag{E91.6}
\]
Here $\mathrm{Art}$ is restriction to the actual ninth-genus
field $\mathcal G_9$ of E.76, not an abstract quotient of the
same order. In its coordinates
$\Gamma_9=C_9^3\oplus C_3$, the actual Kummer character and
mutually inverse kernel coordinates are
\[
 \begin{aligned}
 \chi_9(x,y,z,t)&=8x+7y+2z+3t\pmod9,\\
 A_9\longrightarrow\ker\chi_9:\quad
 (y,z,t)&\longmapsto(-2y+2z+3t,y,z,t),\\
 (x,y,z,t)&\longmapsto(y,z,t).
 \end{aligned}
 \tag{E91.7}
\]
All cyclic moduli are those in the preceding display. In these
$A_9$ coordinates the action, primes and S-quotient are
\[
 \begin{gathered}
 H=\begin{pmatrix}4&3&0\\8&7&6\\1&0&1\end{pmatrix},\qquad
 \mathrm{Art}[\mathfrak p_3]=(1,8,2),\\
 \begin{array}{c|ccc|ccc}
        &Q_1&Q_2&Q_3&T_1&T_2&T_3\\ \hline
 y&6&6&6&6&3&0\\
 z&0&6&3&5&8&5\\
 t&2&2&2&1&1&1
 \end{array},\qquad
 \mathrm{Art}(R_3)=3C_9\oplus C_9\oplus C_3,\\
 A_9/\mathrm{Art}(R_3)\longrightarrow C_3:
       (y,z,t)\longmapsto y\pmod3.
 \end{gathered}
 \tag{E91.8}
\]
The prime labels are transported from E.81 by their unique
totally ramified extensions; $h$ cycles both displayed triples.

These actual action maps impose further unconditional bounds:
\[
 \begin{gathered}
 \exp(\mathcal C_3)=9,\qquad
 \mathcal C_3\in\{C_9^2\oplus C_3, C_9^3, C_9^2\oplus C_3^2\},\\
 [Q_1]+[Q_2]+[Q_3]=3[Q_i]\in W\quad(i=1,2,3),\\
 \operatorname{ord}([Q_2]-[Q_1])=3,\qquad
 [Q_3]-[Q_1]=-([Q_2]-[Q_1]),\\
 \langle[T_1],[T_2],[T_3]\rangle\simeq C_9\oplus C_3,
 \quad \mathrm{Art}\text{ is injective on this subgroup}.
 \end{gathered}
 \tag{E91.9}
\]
The membership in the displayed finite set is a proved group-type
bound. It does not select its remaining alternative.

The exact comparison with the original prime is
\[
 \begin{gathered}
 N_{B_3/L_3}[Q_i]=q_3,\qquad
 j_{B_3/L_3}(q_3)=3[Q_i],\\
 N_{B_3/L_3}(\mathcal C_3)
   =\ker\bigl(\operatorname{Art}_{B_3/L_3}:D_3\to C_3\bigr),
 \qquad [D_3:N(\mathcal C_3)]=3 .
 \end{gathered}
 \tag{E91.10}
\]
In the S-coordinates of E86.4 the last Artin character reads the
first coordinate. In particular the exponent-nine bound for its
norm image remains compatible with both original possibilities
$D_3=C_9\oplus C_3$ and $D_3=C_{27}\oplus C_3$.
The actual value of $3[Q_i]\in W$, and therefore the proposed
all-level principality of the three-prime product, remain a
specified arithmetic step.

One additional computation records the precise result of an
attempted unit construction. Take the full B1 S-unit basis of
E.79, its six valuation rows $V$, and its six cubic local norm
characters $M$. The class matrix of the six B2 S-primes in
their exponent-three subgroup is
\[
 C_S=\begin{pmatrix}
 2&2&2&2&0&1\\1&1&1&2&2&2\\0&0&0&2&2&2
 \end{pmatrix}.
 \quad
 \dim\ker(C_SV\bmod3)=13,
 \quad \operatorname{rank}_{\mathbb F_3}
          M\ker(C_SV)=3.
 \tag{E91.11}
\]
The proof constructs the corresponding actual ordinary B2 units
and compares their B3/B2 symbols. This image equals the already
computed rank-three image of the included ordinary B1 units.
Thus this particular construction has an evaluated image; the
image of all ordinary B2 units is not assigned this rank.

The manuscript anchor remains LM~\S51.16, folios733v--736v,
scans172--178, via the literal positional field norms of E.46,
E.69, E.72 and the all-level divisor comparison of E.88.
\end{statement}

\begin{proof}
\emph{The actual automorphisms and prime restrictions.}
The substitutions on $B_{d+1}$ are
$g_d(\xi)=\xi$, $g_d(\beta_{d+1})=\xi\beta_{d+1}$ and
$h_d(\xi)=\zeta\xi$, $h_d(\beta_{d+1})=\beta_{d+1}$.
They preserve the literal defining relations, have orders nine
and three, and satisfy
\[
 g_d^3=\sigma_d,\qquad
 h_dg_dh_d^{-1}=g_d^4,\qquad
 g_d|_{B_d}=\sigma_{d-1}\quad(d\ge2),\qquad
 h_d|_{B_d}=h_{d-1}.
 \tag{E91.12}
\]
For example the first restriction sends $\beta_d$ to
$\xi^3\beta_d=\zeta\beta_d$ and fixes $\xi$.
These are maps of the indicated fields, with no claim that the
$g_d$ themselves form a compatible sequence.

Every S-prime is totally ramified at every adjacent vertical
step. Hence $\sigma_{d-1}$ fixes each of the six S-primes in
$B_d$, and $g_d$ fixes its unique upper prime in $B_{d+1}$.
The $h_d$ action permutes the three primes over each original
prime in $L_{d+1}$, since $B_{d+1}/L_{d+1}$ is the actual
S-split unramified cubic extension. It is the same permutation
as at B1, transported by the unique vertical prime extension.

For $d\ge2$ put $C=\mathcal C_{d+1}$, $R=R_{d+1}$ and $n=g_d-1$.
E.86 gives $C/R=C_3$. An automorphism of three-power order
acts identically on $C_3$, so $nC\subset R$. The preceding
prime calculation gives $nR=0$. Since $3C\subset R$, it also
gives $3nC=0$. Thus $n^2=0$ and
\[
               (\sigma_d-1)C=((1+n)^3-1)C=0.
 \tag{E91.13}
\]
For $d=1$ the exact action on $\mathcal C_2$ was proved in E86.18.

For either adjacent B or L extension, ordinary ideal norm is
surjective. Here is the class-field map underlying this fact.
If $E/F$ is the indicated cyclic cubic extension and $H/F$
the Hilbert three-class field, then $E\cap H=F$: a nontrivial
intersection would be all of $E$, contradicting its total
ramification at an S-prime. Restriction of Frobenius identifies
the class norm image with
$\operatorname{Gal}(H/(E\cap H))$; equivalently, in the
compositum $EH$ every Frobenius $(a,1)$ supplies a degree-one
prime of $E$ whose ideal norm represents $a$.
This is the actual norm/Artin comparison of class field theory
\cite[V, Proposition3.3, Theorem3.6 and Example3.9;
VIII, Theorem7.4]{E_MilneCFT}.
The identity $jN=\sum\sigma^i$ follows on each prime ideal
from its ramification and residue degrees, and hence on all
classes. Equation (E91.13) gives $jN=3$. It kills the norm
kernel by three. Finally $j(x)=0$ and $N(y)=x$ are equivalent
to $3y=0$, proving the last equality of (E91.2).

\emph{The full abelian-group alternatives for the L tower.}
The ideal $(\beta_d)$ is the unique prime over seventeen.
Thus removing S quotients $D_d$ by precisely $\langle q_d\rangle$.
E.86 identifies that quotient with $C_3^2$, and E.72 gives the
exact order $3^{\alpha_d}$ of its kernel.

For completeness, let a finite abelian three-group $D$ have a
cyclic subgroup $Q=\langle q\rangle$ of order $3^a$, with
$D/Q=C_3^2$ and $a\ge1$. Then $3D\subset Q$, so $3D$ is
cyclic. In the elementary-divisor decomposition at most one
factor can have order greater than three. Counting the order
$|D|=3^{a+2}$, and checking whether $q$ is primitive modulo
$3D$, leaves exactly the two possibilities
\[
 \begin{array}{c|c}
 D&Q\\ \hline
 C_{3^{a+1}}\oplus C_3&3D,\\
 C_{3^a}\oplus C_3^2&\text{a primitive cyclic direct summand}.
 \end{array}
 \tag{E91.14}
\]
To verify the second row directly, a generator $q$ of maximal
order with nonzero reduction can be made the generator of the
long cyclic factor by an automorphism of the displayed group;
its remaining two generators have order three. In the first
row both $Q$ and $3D$ have order $3^a$, giving their equality.
This also treats $a=1$, where the second group is $C_3^3$.

The first row holds at $d=2$ by E.86. Suppose it holds at the
preceding level. In the second alternative at the next level,
the group would be generated by $q_d$ and elements of order
three. Its actual norm, which sends $q_d$ to $q_{d-1}$, could
therefore only have image in
\[
 \langle q_{d-1}\rangle+D_{d-1}[3]
             =3C_{3^{\alpha_{d-1}+1}}\oplus C_3.
 \tag{E91.15}
\]
This has index three in $D_{d-1}$, contradicting the proved
norm surjectivity. Induction proves (E91.3). Its two specified
ideal classes generate because their images form a basis of
$D_d/3D_d$; the elementary finite Nakayama argument lifts that
basis to generators of $D_d$.
E.69 proves $\alpha_d=h$ for all sufficiently large $d$.
The surjective norms then connect finite groups of the same
order and are isomorphisms. Projection onto any such stable
level proves (E91.4), including the actual compatible prime.

\emph{The exact genus quotient and its one possible kernel.}
On $B_3$ the preceding calculation gives $n=g-1$ with
$nR_3=0$, $3n\mathcal C_3=0$ and $n^2=0$. The prime
$\mathfrak p_3$ generates $\mathcal C_3/R_3$ by E.86, so it generates
the image of $n$. The actions of $g$ and $h$ commute on
$\mathcal C_3$, since their field commutator is $g^3=\sigma_2$, whose
class action is trivial. Also $(h-1)\mathcal C_3\subset R_3$. Therefore
$h(n[\mathfrak p_3])=n[\mathfrak p_3]$, giving every assertion
about $W$ in (E91.6).

Let $\mathcal H$ be the full Hilbert three-class field of $B_3$.
It is Galois over $B_1$, since $B_3/B_1$ is cyclic of degree
nine. The commutator subgroup of
$\operatorname{Gal}(\mathcal H/B_1)$ is exactly $(g-1)\mathcal C_3$:
the subgroup $\mathcal C_3$ is abelian and the quotient is cyclic, so
commutators are precisely those with a lift of its generator.
Consequently its maximal abelian-over-$B_1$ subfield is the
genus field $\mathcal G_9$ computed in E.76. Restriction proves
the exact sequence (E91.6), including its kernel.

The retained computation re-evaluates the Kummer character on
the actual ray ideal generators. For a prime coprime to S it
is determined by
$\overline\beta_1^{(N\mathfrak p-1)/9}
 =\overline\xi^{\chi_9(\mathfrak p)}$.
Transport through the full integral Smith maps gives
$(8,7,2,3)$ in the four nontrivial genus coordinates. Substitution
proves both inverse arrows of (E91.7). The actual cyclotomic
automorphism $h$ on the ray quotient gives the matrix $H$.
The six ramified Frobenii are represented by the explicit
Chinese-remainder local uniformizers and coprime ideals of
E.76; their columns are (E91.8). Their lattice image is exactly
$3\mathbb Z\oplus\mathbb Z\oplus\mathbb Z$ modulo
$(9,9,3)$. The prime with $\xi=55$, $\beta_1=22$ has ray
residue $a-\xi=33$ modulo $73$ and gives $(1,8,2)$.
The literal equalities $22^3=63$, $63^3=22$ modulo $73$
identify its norm with the specified $\mathfrak p_3$.
This proves the final S-quotient and every displayed prime map.

\emph{Retaining every central extension and action parameter.}
If $W=0$, the remaining assertions about exponent and primes
follow directly from (E91.8). For the other case let $w$ denote
its actual generator $(g-1)[\mathfrak p_3]_3$, and choose any
lifts $Y,Z,T$ of the three coordinates of $A_9$. The entire
abelian extension has the presentation
\[
 3w=0,\qquad 9Y=a w,\qquad 9Z=b w,\qquad 3T=c w,
             \quad a,b,c\in\mathbb F_3.
 \tag{E91.16}
\]
Indeed these relations give a group of order $3|A_9|$, and
the surjection to $\mathcal C_3$ is therefore an isomorphism. Since
$h(w)=w$, its most general lift of the actual matrix is
\[
 \begin{aligned}
 hY&=4Y+8Z+T+e w,\\
 hZ&=3Y+7Z+f w,\\
 hT&=6Z+T+k w,\qquad e,f,k\in\mathbb F_3.
 \end{aligned}
 \tag{E91.17}
\]
Applying $h$ to $9Y=a w$ forces $b=0$. Put
$A=\operatorname{diag}(9,9,3)$. The exact matrix identity
\[
 A^{-1}(H^3-I)=
 \begin{pmatrix}49&39&24\\112&88&54\\15&12&6\end{pmatrix}
 \tag{E91.18}
\]
then shows that $h^3=1$ forces $a=0$. The correction from
$(e,f,k)$ is zero, because $I+H+H^2$ is zero modulo three.
No value is assigned to $c$. The resulting presentations are
$C_9^2\oplus C_3^2$ for $c=0$ and $C_9^3$ for $c\ne0$.
Together with $W=0$ these prove the exponent and type bounds.

Write $v=(6,0,2)^t$, the actual image of $[Q_1]$. A lift
differs from $6Y+2T$ only by a multiple of $w$. That difference
contributes zero to its $h$-norm. The possible correction from
(E91.17) also vanishes, since $(2I+H)v=0\pmod3$. The exact
relation vectors are
\[
 \begin{aligned}
 A^{-1}(I+H+H^2)v&=(34,80,14)^t,\\
 A^{-1}(3v)&=(2,0,2)^t.
 \end{aligned}
 \tag{E91.19}
\]
With $a=b=0$ both expressions represent $2c w$. Since $h$
cycles the Q-primes, this proves their sum equals three times
each prime class. Their differences are therefore killed by
three; the nonzero genus difference proves the asserted exact
order and the relation between the two differences.

For the T-primes the literal divisor of $\beta_3$ gives
$[T_1]+[T_2]+[T_3]=0$. The matrix image of $[T_1]$ is
$(6,5,1)$, of order nine, while that of $[T_2]-[T_1]$ is
$(6,3,0)$, of order three. In (E91.16), their triples differ
by $-9Y+9Z$ and a multiple of $3w$, which vanish because
$a=b=0$. They consequently generate $C_9\oplus C_3$;
their two genus images already generate a group of order27.
This proves injectivity on the T-subgroup.

\emph{The original prime and the precise limit of the unit attempt.}
The horizontal extension is unramified and each original
S-prime splits into three degree-one primes. Ideal norm and
extension give exactly (E91.10), with the sum now evaluated
by (E91.19). The class norm image has index three, because
$B_3$ itself is the indicated subfield of the Hilbert class
field of $L_3$. Restriction of Frobenius identifies it with
the stated Artin kernel. Thus an exponent-nine norm subgroup
does not give an exponent-nine bound on all of $D_3$; the
actual index-three kernel in $C_{27}\oplus C_3$ can have
exponent nine. The remaining value is stated explicitly.

Finally let $u$ be a B1 S-unit whose valuation vector $v$
has zero class in the B2 S-prime subgroup. The full principal
S-lattice in E86.11 supplies an actual S-unit $w_0$ in B2
with divisor $v$ in the transported prime basis. The included
element $u$ has divisor $3v$, so
$\varepsilon=u/w_0^3$ is an ordinary B2 unit. Its cubic norm
characters for B3/B2 are exactly those of $u$ for B2/B1.
At each of the six completions this is the restriction/norm
identity for the Hilbert pairing,
\[
 (\operatorname{res}u,\beta_2)_{B_2,3}
       =(u,N_{B_2/B_1}\beta_2)_{B_1,3}
       =(u,\beta_1)_{B_1,3}.
 \tag{E91.20}
\]
The cube $w_0^3$ contributes zero. This is the exact pairing
square proved from the restriction/corestriction identity in
E.79 and \cite[III,Theorem3.1]{E_MilneCFT}.
Thus the full image of this construction is
$M\ker(C_SV\bmod3)$. The complete fifteen-column matrices,
the thirteen-dimensional kernel, and its image are retained.
Their ranks are respectively the values in (E91.11), with
the same rank-three image as the first nine ordinary-unit
columns. This records what the attempt computes without
assigning its image to the full ordinary-unit group.

The portable program reconstructs every genus input from the
certified degree-eighteen B1 field and retains full Smith maps
and inverses. It verifies all $243$ simple-coordinate round
trips, all $27$ abelian extension parameters and all $81$
remaining lifts of $h$, as well as the exact class-type
alternatives and the unit-lifting matrix. The initial row-shape
comparison error and its actual output are retained separately;
no assertion uses that failed execution.
\end{proof}
